\documentclass[11pt,a4paper]{article}
\usepackage[utf8]{inputenc}
\usepackage[T1]{fontenc}
\usepackage[spanish,es-tabla]{babel}
\usepackage[a4paper,top=2cm,bottom=2cm,left=2cm,right=2cm]{geometry}
\usepackage{tabularx,booktabs,array,enumitem,amssymb}
\sloppy

\begin{document}

\begin{center}
{\large\textbf{UNIVERSIDAD T\'ECNICA ESTATAL DE QUEVEDO}}\\
Facultad de Ciencias de la Computaci\'on --- Carrera de Ingenier\'ia de Software\\[0.3cm]
\textbf{\Large FORMULARIO DE CONSENTIMIENTO INFORMADO}\\
\textbf{Modelo sin firmar --- no constituye evidencia}
\end{center}

\noindent\begin{tabularx}{\linewidth}{|p{2.8cm}|X|}
\hline
\textbf{C\'odigo} & SGR-IA-BAL-A01 \hfill \textbf{Versi\'on} 1.0\\
\hline
\textbf{Fecha del modelo} & 19 de noviembre de 2026 \hfill \textbf{Categor\'ia} C\\
\hline
\textbf{Asignatura} & Ingenier\'ia de Requerimientos (ISR-401), 4.\textsuperscript{o} nivel, Paralelo B, SPA 2026--2027\\
\hline
\textbf{Proyecto} & Sistema de Gesti\'on para Restaurante con M\'odulo de Inteligencia Artificial (SGR-IA)\\
\hline
\textbf{Organizaci\'on} & El Balcón Restaurant, Quevedo, Los R\'ios\\
\hline
\textbf{Docente} & Gleiston Cicer\'on Guerrero Ulloa, PhD --- gguerrero@uteq.edu.ec\\
\hline
\textbf{Contacto del equipo} & Robinson Espinoza Jeanpierre (l\'ider)\\
\hline
\textbf{Repositorio} & https://github.com/Roselyn15/ISR401-NRS\\
\hline
\textbf{C\'odigo del participante} & \rule{6cm}{0.4pt}\\
\hline
\end{tabularx}

\section*{1. Para qu\'e es este estudio}
Un equipo de estudiantes de Ingenier\'ia de Software caracteriza la operaci\'on de El Balcón Restaurant para delimitar los requisitos de un sistema de gesti\'on asistido por inteligencia artificial. La informaci\'on se usa solo en documentos acad\'emicos de la asignatura. La participaci\'on no implica compra, instalaci\'on ni entrega de un sistema a la organizaci\'on.

\section*{2. Qu\'e tendr\'ia que hacer}
Seg\'un el rol, se solicitar\'a una o varias de estas actividades, indicadas antes de empezar:
\begin{itemize}[leftmargin=16pt,itemsep=1pt]
\item Una entrevista de 15 a 30 minutos sobre tareas diarias y mejoras consideradas.
\item Responder un cuestionario breve, impreso o en l\'inea.
\item Permitir la observaci\'on no participante del trabajo.
\item Revisar contextos y modelos propuestos en una sesi\'on de validaci\'on.
\end{itemize}
Se registran notas an\'onimas. El audio, el video y las fotograf\'ias solo se usan si se autoriza expresamente en la secci\'on~7.

\section*{3. Riesgos y c\'omo se reducen}
No se prev\'en riesgos f\'isicos. Pueden aparecer incomodidad al hablar de problemas internos o uso indebido de la informaci\'on. Se reducen permitiendo omitir preguntas, pausar o terminar la actividad, usando un c\'odigo en lugar del nombre y evitando que una opini\'n identificable llegue a superiores.

\section*{4. Beneficios}
No hay pago ni beneficio econ\'omico. La colaboraci\'on ayuda a que la propuesta responda a necesidades reales y contribuye a la formaci\'on profesional.

\section*{5. Confidencialidad y protecci\'on de datos}
El tratamiento se rige por la Ley Org\'anica de Protecci\'on de Datos Personales del Ecuador (Asamblea Nacional, 2021) y por el C\'odigo de \'Etica y Conducta Profesional de la ACM (2018), seg\'un se referencia en el informe principal.
\begin{itemize}[leftmargin=16pt,itemsep=1pt]
\item \textbf{Nunca se publica:} nombre, c\'edula, firma, rostro, voz, contacto, cargo espec\'ifico ni datos de clientes o proveedores. Endocumentos p\'ublicos aparece solo el c\'odigo de rol.
\item \textbf{D\'onde se guarda lo identificable:} el original firmado y las grabaciones se custodian en carpeta institucional privada \texttt{@uteq.edu.ec}, compartida con el docente.
\item \textbf{Qu\'e se publica:} \'unicamente este modelo en blanco y el manifiesto \texttt{MANIFIESTO.md} con nombre de archivo, fecha, duraci\'on y huella SHA-256. No se publica ninguna copia firmada ni enmascarada.
\item \textbf{Tiempo de conservaci\'on:} el material identificable se conserva como m\'aximo 24 meses desde el cierre del per\'iodo y luego se destruye.
\end{itemize}

\section*{6. Participaci\'on voluntaria}
La participaci\'on es libre. Se puede negar a responder o retirarse en cualquier momento, sin consecuencias. El formulario est\'a dirigido a personas de 18 a\~nos o m\'as.

\section*{7. Autorizaciones espec\'ificas}
Marque S\'i o No. Solo la autorizaci\'on (a) es necesaria para participar.

\noindent\begin{tabularx}{\linewidth}{|X|c|c|}
\hline
\textbf{Autorizaci\'on} & \textbf{S\'i} & \textbf{No}\\
\hline
a) Acepto participar en la actividad explicada. & $\square$ & $\square$\\
\hline
b) Autorizo grabaci\'on en audio. & $\square$ & $\square$\\
\hline
c) Autorizo grabaci\'on en video. & $\square$ & $\square$\\
\hline
d) Autorizo fotograf\'ias del entorno sin rostro, documentos legibles ni GPS. & $\square$ & $\square$\\
\hline
e) Autorizo citas textuales an\'onimas. & $\square$ & $\square$\\
\hline
f) Autorizo uso de datos anonimizados en trabajos de la asignatura. & $\square$ & $\square$\\
\hline
g) Autorizo uso de datos anonimizados en publicaciones derivadas. & $\square$ & $\square$\\
\hline
h) Autorizo dep\'osito en repositorio abierto bajo CC BY 4.0. & $\square$ & $\square$\\
\hline
\end{tabularx}

\section*{8. Declaraci\'on y firmas}
Declaro que he le\'ido este formulario, pude preguntar y comprendo el objetivo, los riesgos, los beneficios, el manejo de mis datos y mi derecho a retirarme.

\vspace{0.6cm}
\noindent\begin{tabularx}{\linewidth}{|X|X|}
\hline
Nombres y apellidos del participante & Firma del participante\\
\hline
\vspace{0.8cm} & \\
\hline
N\'umero de c\'edula & Lugar y fecha (dd/mm/aaaa)\\
\hline
\vspace{0.8cm} & \\
\hline
Nombre y firma de quien aplica & Nombre y firma del testigo (opcional)\\
\hline
\end{tabularx}

\vspace{0.4cm}
\textbf{Uso interno:} el original firmado se guarda en la carpeta privada \texttt{@uteq.edu.ec}. En el repositorio p\'ublico solo se mantiene este modelo sin firmar y el registro \texttt{MANIFIESTO.md}. La fecha del consentimiento debe ser igual o anterior a la de la entrevista, encuesta u observaci\'on que ampara.

\end{document}
